% REVISÕES ------------------------------------------------------

\clearpage
\thispagestyle{empty}
\begin{center}
	\textbf{REVISÕES}
\end{center}

% REVISÕES ------------------------------------------------------
\begin{center} 
	\begin{longtable}{lp{2cm}p{11cm}}
%		\caption{Revisão do documento}\\  % Legenda da tabela
		\hline                              % Linha horizontal
		Revisão & Data & Descrição \\       % Cabeçalho da tabela
		\hline                              % Linha horizontal
		\endfirsthead                       % Fim do cabeçalho da primeira página
		\hline                              % Linha horizontal nas páginas seguintes
		Revisão & Data & Descrição \\       % Cabeçalho repetido
		\hline                              % Linha horizontal após o cabeçalho repetido
		\endhead                            % Fim do cabeçalho repetido

		
		% REVISÕES ------------------------------------------------------
		% rev | mes/ano | descricao
		% REVISÕES ------------------------------------------------------
		% rev | mes/ano | descricao
		00 & 07/2013 & versão inicial\\ \hline
		01 & 12/2013 & correções no perfil e PUDs\\ \hline
		02 & 05/2014 & correções na matriz e atualizações dos PUDs e atualização do TCC\\ \hline
		03 & 09/2014 & correção do PUD de Probabilidade e Estatística\\ \hline
		04 & 06/2017 & atualização de laboratórios, inclusão do LPC e LAMSC\\ \hline
		05 & 06/2017 & inclusão da disciplina de Física III (04506.18) como pré-requisito na disciplina de Máquinas Elétricas (04506.36)\\ \hline
		06 & 06/2017 & retirada da disciplina de Algebra Linear (04506.1) como pré-requisito na disciplina de Cálculo II (04506.7)\\ \hline
		07 & 06/2017 & inclusão da disciplina de Equações Diferenciais Ordinárias (04507.60) como optativa\\ \hline
		08 & 06/2017 & inclusão da disciplina de Inteligência Computacional Aplicada (04507.61) como optativa\\ \hline
		09 & 06/2017 & inclusão da disciplina de Sistemas Lineares (04507.26) como pré-requisito na disciplina de Controle I (04507.31)\\ \hline
		10 & 07/2017 & inclusão da disciplina de Educação Física (04507.59) como optativa\\ \hline
		11 & 07/2017 & inclusão da disciplina de Mecânica dos Fluidos (04507.62) como optativa\\ \hline
		12 & 08/2017 & inclusão da disciplina de Eletrônica I (04507.15) como pré-requisito na disciplina de Controladores Lógicos Programáveis (04507.35)\\ \hline
		13 & 08/2017 & alteracao da carga horaria da disciplina de Libras de 40hs para 80hs\\ \hline
		14 & 08/2017 & inclusão da disciplina de Controle I (04507.31) como pré-requisito na disciplina de Identificação de Sistemas (04507.56)\\ \hline
		%14 & 08/2017 & alteracao da carga horaria da disciplina de Projetos Sociais de 40hs para 80hs\\
		15 & 08/2017 & adequação do PPC conforme instrumental da PROEN\\ \hline
		16 & 11/2017 & atualização de PPC aprovada pelo Consup, resolução anterior: 002-14, nova: 117-17\\ \hline
		% REVISÕES DE 2019 ---------------------------------------------------------------
		17 & 05/2019 & mudança do pré-requisito de Circuitos Elétricos I (04507.20), que passou de Cálculo 3 (04507.12) para Cálculo 1 (04507.2) \\ \hline
		18 & 05/2019 & mudança do pré-requisito de Controle I (04507.31), que passou de Sistemas Lineares (04507.26) para Circuitos Elétricos II (04507.24) \\ \hline
		19 & 05/2019 & mudança do pré-requisito de Processamento Digital de Sinais (04507.34), que passou de Linguagem de Programação (04507.14) para Sistemas Lineares (04507.26) \\ \hline
		20 & 05/2019 & mudança do pré-requisito de Controladores Lógicos Programáveis (04507.35), que passou de Eletrônica I (04507.15) para Microcontroladores (04507.28) \\ \hline
		21 & 05/2019 & mudança do pré-requisito de Acionamentos Hidráulicos e Pneumáticos (04507.37), que passou de Eletrônica II (04507.23) para Instalações Elétricas (04507.30) \\ \hline
		22 & 05/2019 & inclusão da disciplina de Controle I (04507.31) como pré-requisito na disciplina de Controle II (04507.40) \\ \hline
		23 & 05/2019 & inclusão da disciplina de Controle I (04507.31) como pré-requisito na disciplina de Acionamento de Máquinas (04507.41) \\ \hline
		24 & 05/2019 & mudança dos pré-requisitos de Robótica II (04507.55), que passaram de Linguagem de Programação (04507.14) e Dispositivos Periféricos (04507.42) para Robótica I (04507.44) \\ \hline
		25 & 05/2019 & inclusão da disciplina de Controle II (04507.40) como pré-requisito na disciplina de Controle III (04507.58) \\ \hline
		26 & 05/2019 & modificação do formato dos PUDs, nos quais foram inseridas células relativas aos seguintes tópicos: Carga Horária Teórica, Carga Horária Prática, Interdisciplinaridade, Temas Atuais e Recursos \\ \hline
		27 & 05/2019 & revisões de diversos PUDs, primeiramente para preencher as novas células inseridas, bem como para readequar os seguintes espaços: Ementa, Objetivos, Programa, Metodologia, Avaliação e Bibliografia \\ \hline
		28 & 05/2019 & revisões de textos ao longo de todo o PPC, como forma a atender orientações da CTP do campus, em especial nos seguintes tópicos: Corpo técnico-administrativo (item 18), Departamento de Extensão, Pesquisa, Pós-Graduação e Inovação (subitem 16.3), Setor de Estágio (subitem 16.8), Colegiado (subitem 12.3), Ações decorrentes dos processos de autoavaliação e avaliação externa (subitem 12.5), Atividades complementares (subitem 10.5), Coordenaçao de curso (subitem 16.1) \\ \hline
		29 & 05/2019 & revisão do texto relativo ao TCC (subitem 10.6), em especial nos seguintes itens: texto relativo à entrega da versão para apresentação pública do TCC; retirada a necessidade de entrega de 01 (uma) cópia impressa da versão final do TCC, encadernada em capa dura com letras douradas, sendo necessária apenas a versão digital; inserido texto que deixa claro que o arquivo eletrônico em mídia digital do TCC será disponibilizado para acesso público no Sistema de Bibliotecas do IFCE - SIBI; inserido o endereço para o Manual de Normalização de Trabalhos Acadêmicos do IFCE; inserido texto que condiciona a matricula do aluno na disciplina de TCC à quantidade de créditos concluídos (80\%) e não ao semestre.\\ \hline
		30 & 05/2021 & atualização das composições do Colegiado e do NDE do curso, atualização do texto relativo ao estágio curricular (subitem 10.4), permitindo a realização dos estágios obrigatórios e não obrigatórios de forma remota; atualização do corpo docente do curso (subitens 17.2 e 17.3), em virtude dos recentes processos de remoção.\\ \hline
		31 & 09/2021 & Inserção da Portaria 430 do Ministério da Educação, de 3 de maio de 2021, que trata do reconhecimento do curso de Engenharia de Controle e Automação do IFCE - Campus Maracanaú; atualização da composição do Colegiado.\\ \hline
		32 & 12/2025 & {
			\vm 
			Reformulação do PPC do curso para atender às instruções: a) Resolução CNE/CES nº 2, de 24 de abril de 2019;
			b) Resolução CNE/CES nº 1, de 26 de março de 2021; c) Resolução CONSUP/IFCE Nº 259, de 08 de Janeiro de 2025 (Manual de Normalização de Projetos Pedagógicos dos Cursos de Engenharia do IFCE); d) Guia de curricularização das atividades de extensão nos cursos técnicos, de graduação e pós-graduação do IFCE; e) Instrumento de avaliação de cursos de graduação, presencial e a distância, voltado ao reconhecimento e à renovação de reconhecimento, conforme o Sistema Nacional de Avaliação da Educação Superior.
		} \\ \hline
		
		
		% REVISÕES ------------------------------------------------------
        
	\end{longtable}
\end{center}
